% GENERATED FILE. Edit canonical lessons, then run npm run generate:books.
% canonical-source-hash: fnv1a64:2a66dd96999f1a01
% canonical-lessons: TE-C534-godavapadu, TE-C534-tagada, TE-C534-alugu, TE-C534-raji, TE-C534-manaspartha

\chapter{Quarrels}
\label{ch:te-quarrels}
\hlchaptermodality{534}

\begin{chapteropening}
\textbf{By the end of this chapter:} \emph{I can talk about rows, disputes, sulks and making up.}

Complete the last lesson of chapter 534: I can talk about rows, disputes, sulks and making up.
\end{chapteropening}

\section[\texorpdfstring{\te{గొడవపడు} (goḍavapaḍu)}{గొడవపడు (goḍavapaḍu)}]{\te{గొడవపడు} (goḍavapaḍu) — to have a row; to quarrel}

\label{lesson:TE-C534-godavapadu}



\begin{quote}
Before the new one: say the Telugu for hospitality, then the Telugu for to chat.
\end{quote}

\subsection*{You'll want to know: \te{గొడవపడు}}
\textbf{\te{గొడవపడు}} — \emph{goḍavapaḍu} — \textquotedblleft{}to have a row; to quarrel\textquotedblright{}.

\begin{quote}
\te{చిన్న} \te{విషయానికి} \te{అన్నదమ్ములు} \te{గొడవపడ్డారు}.
\end{quote}

\begin{quote}
\emph{cinna viṣayāniki annadammulu goḍavapaḍḍāru.}
\end{quote}

\begin{quote}
The brothers had a row over a small thing.
\end{quote}

Past (I …): \te{గొడవపడ్డాను}, \emph{goḍavapaḍḍānu}.

-\te{తో} \te{గొడవపడు}, to quarrel with.

\begin{cousinweb}
\textbf{\te{గొడవపడు}} (\emph{goḍavapaḍu}) is native Dravidian: DEDR 2048 lists Telugu \emph{goḍava}, \textquotedblleft{}bother, worry, fuss\textquotedblright{}, with Kannada \emph{goḍave}; the native verb \emph{paḍu}, \textquotedblleft{}to feel, to suffer\textquotedblright{} (DEDR 3853) makes the verb.
\end{cousinweb}

\begin{grammarlens}[title={What you've built}]
One more word for describing what you lived and felt.
\end{grammarlens}

\subsection*{Your turn}
\begin{itemize}
\raggedright
  \item \emph{Say it:} \emph{goḍavapaḍu}
  \item \emph{Say it:} \emph{goḍavapaḍu}, once more
  \item \emph{Say it:} \emph{cinna viṣayāniki annadammulu goḍavapaḍḍāru.} — The brothers had a row over a small thing.
  \item \emph{Recall:} say the Telugu for hospitality, then the Telugu for to chat, then say \emph{goḍavapaḍu} again
\end{itemize}

\begin{tcolorbox}[breakable,colback=teal!4,colframe=teal!35!black,title={Before you move on}]
What does \te{గొడవపడు} mean? (\textquotedblleft{}To have a row; to quarrel\textquotedblright{}.) Say it once more.
\end{tcolorbox}
\section[\texorpdfstring{\te{తగాదా} (tagādā)}{తగాదా (tagādā)}]{\te{తగాదా} (tagādā) — a dispute; a row}

\label{lesson:TE-C534-tagada}



\begin{quote}
Before the new one: say the Telugu for to chat, then the Telugu for to have a row.
\end{quote}

\subsection*{You'll want to know: \te{తగాదా}}
\textbf{\te{తగాదా}} — \emph{tagādā} — \textquotedblleft{}a dispute; a row\textquotedblright{}.

\begin{quote}
\te{పొలం} \te{గురించి} \te{రెండు} \te{కుటుంబాల} \te{మధ్య} \te{తగాదా} \te{వచ్చింది}.
\end{quote}

\begin{quote}
\emph{polaṁ guriñci reṇḍu kuṭumbāla madhya tagādā vaccindi.}
\end{quote}

\begin{quote}
A dispute arose between two families over a field.
\end{quote}

Often about land or money.

\begin{cousinweb}
\textbf{\te{తగాదా}} (\emph{tagādā}) is usually traced to Arabic \emph{taqāḍā}, \textquotedblleft{}a demand, a claim\textquotedblright{}, which came through Hindustani.
\end{cousinweb}

\begin{grammarlens}[title={What you've built}]
One more word for describing what you lived and felt.
\end{grammarlens}

\subsection*{Your turn}
\begin{itemize}
\raggedright
  \item \emph{Say it:} \emph{tagādā}
  \item \emph{Say it:} \emph{tagādā}, once more
  \item \emph{Say it:} \emph{polaṁ guriñci reṇḍu kuṭumbāla madhya tagādā vaccindi.} — A dispute arose between two families over a field.
  \item \emph{Recall:} say the Telugu for to chat, then the Telugu for to have a row, then say \emph{tagādā} again
\end{itemize}

\begin{tcolorbox}[breakable,colback=teal!4,colframe=teal!35!black,title={Before you move on}]
What does \te{తగాదా} mean? (\textquotedblleft{}A dispute; a row\textquotedblright{}.) Say it once more.
\end{tcolorbox}
\section[\texorpdfstring{\te{అలుగు} (alugu)}{అలుగు (alugu)}]{\te{అలుగు} (alugu) — to sulk; to take offence}

\label{lesson:TE-C534-alugu}



\begin{quote}
Before the new one: say the Telugu for to have a row, then the Telugu for a dispute.
\end{quote}

\subsection*{You'll want to know: \te{అలుగు}}
\textbf{\te{అలుగు}} — \emph{alugu} — \textquotedblleft{}to sulk; to take offence\textquotedblright{}.

\begin{quote}
\te{బొమ్మ} \te{కొనలేదని} \te{పాప} \te{అలిగింది}.
\end{quote}

\begin{quote}
\emph{bomma konalēdani pāpa aligindi.}
\end{quote}

\begin{quote}
The little girl sulked because nobody bought her the doll.
\end{quote}

Past (I …): \te{అలిగాను}, \emph{aligānu}.

The noun is \te{అలక}, a sulk.

\begin{cousinweb}
\textbf{\te{అలుగు}} (\emph{alugu}) is native Dravidian: DEDR 304 lists Telugu \emph{alugu}, \textquotedblleft{}to be angry or displeased, to take offence\textquotedblright{}, and \emph{aluka}, \textquotedblleft{}displeasure\textquotedblright{}.
\end{cousinweb}

\begin{grammarlens}[title={What you've built}]
One more word for describing what you lived and felt.
\end{grammarlens}

\subsection*{Your turn}
\begin{itemize}
\raggedright
  \item \emph{Say it:} \emph{alugu}
  \item \emph{Say it:} \emph{alugu}, once more
  \item \emph{Say it:} \emph{bomma konalēdani pāpa aligindi.} — The little girl sulked because nobody bought her the doll.
  \item \emph{Recall:} say the Telugu for to have a row, then the Telugu for a dispute, then say \emph{alugu} again
\end{itemize}

\begin{tcolorbox}[breakable,colback=teal!4,colframe=teal!35!black,title={Before you move on}]
What does \te{అలుగు} mean? (\textquotedblleft{}To sulk; to take offence\textquotedblright{}.) Say it once more.
\end{tcolorbox}
\section[\texorpdfstring{\te{రాజీ} (rājī)}{రాజీ (rājī)}]{\te{రాజీ} (rājī) — a compromise; a reconciliation}

\label{lesson:TE-C534-raji}



\begin{quote}
Before the new one: say the Telugu for a dispute, then the Telugu for to sulk.
\end{quote}

\subsection*{You'll want to know: \te{రాజీ}}
\textbf{\te{రాజీ}} — \emph{rājī} — \textquotedblleft{}a compromise; a reconciliation\textquotedblright{}.

\begin{quote}
\te{చివరికి} \te{ఇద్దరూ} \te{రాజీ} \te{పడ్డారు}.
\end{quote}

\begin{quote}
\emph{civariki iddarū rājī paḍḍāru.}
\end{quote}

\begin{quote}
In the end the two of them came to terms.
\end{quote}

\te{రాజీ} \te{పడు}, to come to terms.

\begin{cousinweb}
\textbf{\te{రాజీ}} (\emph{rājī}) came into Telugu through Hindustani \emph{rāzī}, \textquotedblleft{}agreed, content\textquotedblright{}, which is usually traced to Arabic.
\end{cousinweb}

\begin{grammarlens}[title={What you've built}]
One more word for describing what you lived and felt.
\end{grammarlens}

\subsection*{Your turn}
\begin{itemize}
\raggedright
  \item \emph{Say it:} \emph{rājī}
  \item \emph{Say it:} \emph{rājī}, once more
  \item \emph{Say it:} \emph{civariki iddarū rājī paḍḍāru.} — In the end the two of them came to terms.
  \item \emph{Recall:} say the Telugu for a dispute, then the Telugu for to sulk, then say \emph{rājī} again
\end{itemize}

\begin{tcolorbox}[breakable,colback=teal!4,colframe=teal!35!black,title={Before you move on}]
What does \te{రాజీ} mean? (\textquotedblleft{}A compromise; a reconciliation\textquotedblright{}.) Say it once more.
\end{tcolorbox}
\section[\texorpdfstring{\te{మనస్పర్థ} (manaspartha)}{మనస్పర్థ (manaspartha)}]{\te{మనస్పర్థ} (manaspartha) — a rift; a misunderstanding}

\label{lesson:TE-C534-manaspartha}



\begin{quote}
Before the new one: say the Telugu for to sulk, then the Telugu for a compromise.
\end{quote}

\subsection*{You'll want to know: \te{మనస్పర్థ}}
\textbf{\te{మనస్పర్థ}} — \emph{manaspartha} — \textquotedblleft{}a rift; a misunderstanding\textquotedblright{}.

\begin{quote}
\te{చిన్న} \te{మనస్పర్థ} \te{వల్ల} \te{స్నేహితులు} \te{మాట్లాడుకోవడం} \te{మానేశారు}.
\end{quote}

\begin{quote}
\emph{cinna manaspartha valla snēhitulu māṭlāḍukōvaḍaṁ mānēśāru.}
\end{quote}

\begin{quote}
Because of a small rift the friends stopped talking.
\end{quote}

Plural \te{మనస్పర్థలు}; \te{మనసు}, mind, + \te{స్పర్థ}, rivalry.

\begin{cousinweb}
\textbf{\te{మనస్పర్థ}} (\emph{manaspartha}) is a learned compound built from Sanskrit \emph{manas}, \textquotedblleft{}mind\textquotedblright{}, and \emph{spardhā}, \textquotedblleft{}rivalry, emulation\textquotedblright{} (Monier-Williams).
\end{cousinweb}

\begin{grammarlens}[title={What you've built}]
One more word for describing what you lived and felt.
\end{grammarlens}

\subsection*{Your turn}
\begin{itemize}
\raggedright
  \item \emph{Say it:} \emph{manaspartha}
  \item \emph{Say it:} \emph{manaspartha}, once more
  \item \emph{Say it:} \emph{cinna manaspartha valla snēhitulu māṭlāḍukōvaḍaṁ mānēśāru.} — Because of a small rift the friends stopped talking.
  \item \emph{Recall:} say the Telugu for to sulk, then the Telugu for a compromise, then say \emph{manaspartha} again
\end{itemize}

\begin{tcolorbox}[breakable,colback=teal!4,colframe=teal!35!black,title={Before you move on}]
What does \te{మనస్పర్థ} mean? (\textquotedblleft{}A rift; a misunderstanding\textquotedblright{}.) Say it once more.
\end{tcolorbox}
